\chapter{Networks}
%%%%%%%%%%%%%%%%%%%%%%%%%%%%%%%%%%%%%%%%%%%%%%%%%%%%%%%%%%%%%%%%%%%%%%%%%%%%%%%
\section{Networks}

% Generated by scripts/blueprint_restate.py from TCSlib.GraphTheory.Core.Networks.
% Each body is the STATEMENT only, taken from the declaration's Lean docstring:
% the one-line summary plus the verbatim book statement.  Proof, proof plan,
% significance commentary and formalisation notes are excluded by construction.

\begin{definition}[A network $N = (D, X, Y, c)$]
\lean{Networks.Network}
\label{Networks.Network}
\leanok
A \textbf{network} \texttt{N = (D, X, Y, c)}.

\emph{A network \texttt{N} is a digraph \texttt{D} (the
underlying digraph of \texttt{N}) with two distinguished subsets of vertices, \texttt{X}
and \texttt{Y},
and a non-negative integer-valued function \texttt{c} defined on its arc set \texttt{A};
the sets
\texttt{X} and \texttt{Y} are assumed to be disjoint and nonempty.  The vertices in \texttt{X} are the
sources of \texttt{N} and those in \texttt{Y} are the sinks of \texttt{N}. They
correspond to
production centres and markets, respectively. \dots{} The function \texttt{c} is the
capacity
function of \texttt{N} and its value on an arc \texttt{a} the capacity of \texttt{a}.
The capacity of
an arc can be thought of as representing the maximum rate at which a commodity can
be transported along it.}
\end{definition}

\begin{definition}[Intermediate vertices $I = V \ (X \cup Y)$]
\lean{Networks.Network.I}
\label{Networks.Network.I}
\leanok
\uses{Networks.Network}
Intermediate vertices \texttt{I = V \textbackslash{} (X $\cup$ Y)}.

\emph{Vertices which are neither sources nor
sinks are called intermediate vertices; the set of such vertices will be denoted by
\texttt{I}.}
\end{definition}

\begin{definition}[$f^{+}(S) = f(S, \bar{S}) = \sum_{u\in S} \sum_{v\notin S} f u v$]
\lean{Networks.Network.fOut}
\label{Networks.Network.fOut}
\leanok
\uses{Networks.Network}
\texttt{f$^{+}$(S) = f(S, $\bar{S}$) = $\sum$\_\{u$\in$S\} $\sum$\_\{v$\notin$S\} f u v}.

\emph{If \texttt{S $\subseteq$ V}, we denote \texttt{V \textbackslash{} S} by
\texttt{$\bar{S}$}. \dots{}
If \texttt{f} is a real-valued function defined on the arc set \texttt{A} of \texttt{N},
and if
\texttt{K $\subseteq$ A}, we denote \texttt{$\sum$\_\{a $\in$ K\} f(a)} by \texttt{f(K)}.  Furthermore, if \texttt{K} is a set of
arcs of the form \texttt{(S, $\bar{S}$)}, we shall write \texttt{f$^{+}$(S)} for
\texttt{f(S, $\bar{S}$)} and \texttt{f$^{-}$(S)} for
\texttt{f($\bar{S}$, S)}.}
\end{definition}

\begin{definition}[$f^{-}(S) = f(\bar{S}, S) = \sum_{u\notin S} \sum_{v\in S} f u v$]
\lean{Networks.Network.fIn}
\label{Networks.Network.fIn}
\leanok
\uses{Networks.Network}
\texttt{f$^{-}$(S) = f($\bar{S}$, S) = $\sum$\_\{u$\notin$S\} $\sum$\_\{v$\in$S\} f u v}.

\emph{\dots{}and \texttt{f$^{-}$(S)} for \texttt{f($\bar{S}$, S)}.}

The book names the difference: \emph{If \texttt{S} is a subset of vertices in a network
\texttt{N} and
\texttt{f} is a flow in \texttt{N}, then \texttt{f$^{+}$(S) - f$^{-}$(S)} is called the resultant flow out of \texttt{S},
and \texttt{f$^{-}$(S) - f$^{+}$(S)} the resultant flow into \texttt{S}, relative to
\texttt{f}.}
\end{definition}

\begin{definition}[A flow: (11.1) capacity constraint + (11.2) conservation at intermediate vertices]
\lean{Networks.Network.IsFlow}
\label{Networks.Network.IsFlow}
\leanok
\uses{Networks.Network, Networks.Network.I, Networks.Network.fIn, Networks.Network.fOut}
A \textbf{flow}: (11.1) capacity constraint + (11.2) conservation at intermediate
vertices.

\emph{A flow in a network \texttt{N} is an
integer-valued function \texttt{f} defined on \texttt{A} such that}

0 $\le$ f(a) $\le$ c(a)   for all a $\in$ A                                  (11.1)
f$^{-}$(v) = f$^{+}$(v) for all v $\in$ I (11.2)

\emph{The upper bound in condition (11.1) is called the capacity constraint; it imposes
the natural restriction that the rate of flow along an arc cannot exceed the
capacity of the arc.  Condition (11.2), called the conservation condition, requires
that, for any intermediate vertex \texttt{v}, the rate at which material is transported
into \texttt{v} is equal to the rate at which it is transported out of \texttt{v}. Note
that
every network has at least one flow, since the function \texttt{f} defined by
\texttt{f(a) = 0},
for all \texttt{a $\in$ A}, clearly satisfies both (11.1) and (11.2); it is called the
zero
flow.}
\end{definition}

\begin{definition}[\texttt{val f = f$^{+}$(X) $-$ f$^{-}$(X)}]
\lean{Networks.Network.val}
\label{Networks.Network.val}
\leanok
\uses{Networks.Network, Networks.Network.fIn, Networks.Network.fOut}
\texttt{val f = f$^{+}$(X) $-$ f$^{-}$(X)}.

\emph{Since the conservation condition
requires that the resultant flow out of any intermediate vertex is zero, it is
intuitively clear and not difficult to show (exercise 11.1.3) that, relative to any
flow \texttt{f}, the resultant flow out of \texttt{X} is equal to the resultant flow
into \texttt{Y}.
This common quantity is called the value of \texttt{f}, and is denoted by \texttt{val
f}; thus}

val f = f$^{+}$(X) - f$^{-}$(X)

\emph{A flow \texttt{f} in \texttt{N} is a maximum flow if there is no flow \texttt{f'}
in \texttt{N} such that
\texttt{val f' > val f}.}
\end{definition}

\begin{definition}[A cut $(S, \bar{S})$ with $x \in S$, $y \notin S$]
\lean{Networks.Network.IsCut}
\label{Networks.Network.IsCut}
\leanok
\uses{Networks.Network}
A \textbf{cut} \texttt{(S, $\bar{S}$)} with \texttt{x $\in$ S}, \texttt{y $\notin$ S}.

\emph{Let \texttt{N} be a network with a single
source \texttt{x} and a single sink \texttt{y}. A cut in \texttt{N} is a set of arcs of
the form
\texttt{(S, $\bar{S}$)}, where \texttt{x $\in$ S} and \texttt{y $\in$ $\bar{S}$}.}
\end{definition}

\begin{definition}[\texttt{cap (S, $\bar{S}$) = $\sum$\_\{u$\in$S\} $\sum$\_\{v$\notin$S\} c u v}]
\lean{Networks.Network.capOf}
\label{Networks.Network.capOf}
\leanok
\uses{Networks.Network}
\texttt{cap (S, $\bar{S}$) = $\sum$\_\{u$\in$S\} $\sum$\_\{v$\notin$S\} c u v}.

\emph{The capacity of a cut \texttt{K} is the sum of
the capacities of its arcs. We denote the capacity of \texttt{K} by \texttt{cap K}; thus
\texttt{cap K = $\sum$\_\{a $\in$ K\} c(a)}.}
\end{definition}

\begin{definition}[A maximum flow]
\lean{Networks.Network.IsMaxFlow}
\label{Networks.Network.IsMaxFlow}
\leanok
\uses{Networks.Network, Networks.Network.IsFlow, Networks.Network.val}
A \textbf{maximum flow}.

\emph{A flow \texttt{f} in \texttt{N} is a maximum flow if
there is no flow \texttt{f'} in \texttt{N} such that \texttt{val f' > val f}.}
\end{definition}

\begin{definition}[A minimum cut]
\lean{Networks.Network.IsMinCut}
\label{Networks.Network.IsMinCut}
\leanok
\uses{Networks.Network, Networks.Network.IsCut, Networks.Network.capOf}
A \textbf{minimum cut}.

\emph{A cut \texttt{K} in \texttt{N} is a minimum cut if
there is no cut \texttt{K'} in \texttt{N} such that \texttt{cap K' < cap K}.}
\end{definition}

\begin{definition}[An $f$-incrementing path from the source]
\lean{Networks.Network.IncPath}
\label{Networks.Network.IncPath}
\leanok
\uses{Networks.Network}
An \texttt{f}-\textbf{incrementing path} from the source.

\emph{The path \texttt{P} is \texttt{f}-unsaturated if \texttt{$\iota$(P) > 0} (or,
equivalently, if each forward arc of \texttt{P} is \texttt{f}-unsaturated and each
reverse arc of
\texttt{P} is \texttt{f}-positive).  Put simply, an \texttt{f}-unsaturated path is one that is not being
used to its full capacity. An \texttt{f}-incrementing path is an \texttt{f}-unsaturated
path
from the source \texttt{x} to the sink \texttt{y}.}
\end{definition}

\begin{definition}[The residual capacity $\iota(P) > 0$ of an incrementing path --- a recursion over \texttt{IncPath}]
\lean{Networks.Network.IncPath.iota}
\label{Networks.Network.IncPath.iota}
\uses{Networks.Network, Networks.Network.IncPath}
The residual capacity \texttt{$\iota$(P) > 0} of an incrementing path --- a recursion
over \texttt{IncPath}.

\emph{With each path \texttt{P} in \texttt{N} we associate a
non-negative integer \texttt{$\iota$(P) = min\_\{a $\in$ A(P)\} $\iota$(a)}, where
\texttt{$\iota$(a) = c(a) - f(a)} if
\texttt{a} is a forward arc of \texttt{P}, and \texttt{$\iota$(a) = f(a)} if \texttt{a} is a reverse arc.  As may
easily be seen, \texttt{$\iota$(P)} is the largest amount by which the flow along
\texttt{P} can be
increased (relative to \texttt{f}) without violating condition (11.1).}
\end{definition}

\begin{definition}[The revised flow $\hat{f}$ (11.9) obtained by pushing $\iota(P)$ along $P$]
\lean{Networks.Network.revisedFlow}
\label{Networks.Network.revisedFlow}
\uses{Networks.Network, Networks.Network.IncPath}
The revised flow \texttt{$\hat{f}$} (11.9) obtained by pushing \texttt{$\iota$(P)} along
\texttt{P}.

\emph{By sending an additional flow of \texttt{$\iota$(P)} along
\texttt{P}, one obtains a new flow \texttt{$\hat{f}$} defined by \texttt{$\hat{f}$(a) = f(a) + $\iota$(P)} if \texttt{a} is a
forward arc of \texttt{P}, \texttt{$\hat{f}$(a) = f(a) - $\iota$(P)} if \texttt{a} is a
reverse arc of \texttt{P}, and
\texttt{$\hat{f}$(a) = f(a)} otherwise.}
\end{definition}

\begin{definition}[Edge connectivity $\kappa'(G)$]
\lean{Networks.edgeConnectivity}
\label{Networks.edgeConnectivity}
\leanok
\emph{We \dots{} define the edge connectivity
\texttt{$\kappa$'(G)} of \texttt{G} to be the minimum \texttt{k} for which \texttt{G} has a \texttt{k}-edge cut. \dots{}  \texttt{G} is
said to be \texttt{k}-edge-connected if \texttt{$\kappa$'(G) $\ge$ k}.}
\end{definition}

\begin{definition}[Vertex connectivity $\kappa(G)$]
\lean{Networks.vertexConnectivity}
\label{Networks.vertexConnectivity}
\leanok
\emph{If \texttt{G} has at least one pair of distinct
nonadjacent vertices, the connectivity \texttt{$\kappa$(G)} of \texttt{G} is the minimum
\texttt{k} for which
\texttt{G} has a \texttt{k}-vertex cut; otherwise, we define \texttt{$\kappa$(G)} to be \texttt{$\nu$ - 1}. \dots{}  \texttt{G} is said
to be \texttt{k}-connected if \texttt{$\kappa$(G) $\ge$ k}.}
\end{definition}

\begin{definition}[Internally disjoint family of paths]
\lean{Networks.InternallyDisjoint}
\label{Networks.InternallyDisjoint}
\leanok
A family of paths is internally disjoint when no
vertex of \texttt{G} is an internal vertex of more than one path of the family.
\end{definition}

\begin{definition}[Max number of arc-disjoint directed $(x,y)$-paths in a unit-ish \texttt{Network}]
\lean{Networks.Network.maxArcDisjointPaths}
\label{Networks.Network.maxArcDisjointPaths}
\uses{Networks.Network}
Max number of arc-disjoint directed \texttt{(x,y)}-paths in a unit-ish \texttt{Network}.

(B\&M lemma 11.4(a), verbatim). \emph{\dots{}the maximum number \texttt{m} of
arc-disjoint directed \texttt{(x, y)}-paths in \texttt{N}.}
\end{definition}

\begin{definition}[Max number of arc-disjoint directed $(x,y)$-paths in a unit-ish \texttt{Network}]
\lean{Networks.Network.minArcsDestroyingPaths}
\label{Networks.Network.minArcsDestroyingPaths}
\uses{Networks.Network}
Min number of arcs whose deletion destroys all directed \texttt{(x,y)}-paths in a
\texttt{Network}.

(B\&M lemma 11.4(b), verbatim). \emph{\dots{}the minimum number \texttt{n} of arcs whose
deletion destroys all directed \texttt{(x, y)}-paths in \texttt{N}.}
\end{definition}

\begin{definition}[Max number of arc-disjoint directed $(x,y)$-paths in a digraph]
\lean{Networks.maxArcDisjointDirectedPaths}
\label{Networks.maxArcDisjointDirectedPaths}
Max number of arc-disjoint directed \texttt{(x,y)}-paths in a digraph.
\end{definition}

\begin{definition}[Max number of arc-disjoint directed $(x,y)$-paths in a digraph]
\lean{Networks.minArcsDestroyingDirectedPaths}
\label{Networks.minArcsDestroyingDirectedPaths}
Min number of arcs destroying all directed \texttt{(x,y)}-paths in a digraph.
\end{definition}

\begin{definition}[Min number of arcs destroying all directed $(x,y)$-paths in a digraph]
\lean{Networks.maxInternallyDisjointDirectedPaths}
\label{Networks.maxInternallyDisjointDirectedPaths}
Max number of internally-disjoint directed \texttt{(x,y)}-paths in a digraph.
\end{definition}

\begin{definition}[Max number of internally-disjoint directed $(x,y)$-paths in a digraph]
\lean{Networks.minVerticesDestroyingDirectedPaths}
\label{Networks.minVerticesDestroyingDirectedPaths}
Min number of vertices destroying all directed \texttt{(x,y)}-paths in a digraph.
\end{definition}

\begin{definition}[Max number of edge-disjoint $(x,y)$-paths in a graph]
\lean{Networks.maxEdgeDisjointPaths}
\label{Networks.maxEdgeDisjointPaths}
Max number of edge-disjoint \texttt{(x,y)}-paths in a graph.
\end{definition}

\begin{definition}[Max number of edge-disjoint $(x,y)$-paths in a graph]
\lean{Networks.minEdgesDestroyingPaths}
\label{Networks.minEdgesDestroyingPaths}
Min number of edges destroying all \texttt{(x,y)}-paths in a graph.
\end{definition}

\begin{definition}[Min number of edges destroying all $(x,y)$-paths in a graph]
\lean{Networks.maxInternallyDisjointPaths}
\label{Networks.maxInternallyDisjointPaths}
Max number of internally-disjoint \texttt{(x,y)}-paths in a graph.
\end{definition}

\begin{definition}[Max number of internally-disjoint $(x,y)$-paths in a graph]
\lean{Networks.minVerticesDestroyingPaths}
\label{Networks.minVerticesDestroyingPaths}
Min number of vertices destroying all \texttt{(x,y)}-paths in a graph.
\end{definition}

\begin{definition}[Max number of vertex-disjoint $S$--$T$ paths (Ex 11.4.3)]
\lean{Networks.maxSTVertexDisjointPaths}
\label{Networks.maxSTVertexDisjointPaths}
Max number of vertex-disjoint \texttt{S}--\texttt{T} paths (Ex 11.4.3).

\emph{\dots{}the maximum number of
vertex-disjoint paths with one end in \texttt{S} and one end in \texttt{T}.}
\end{definition}

\begin{definition}[Max number of vertex-disjoint $S$--$T$ paths (Ex 11.4.3)]
\lean{Networks.minSTSeparator}
\label{Networks.minSTSeparator}
Min number of vertices separating \texttt{S} from \texttt{T} (Ex 11.4.3).

\emph{\dots{}the minimum number of vertices
whose deletion separates \texttt{S} from \texttt{T} (that is, after deletion no
component contains
a vertex of \texttt{S} and a vertex of \texttt{T}).}
\end{definition}

\begin{definition}[The associated digraph $D(G)$: each edge becomes two opposite arcs (cited as ex 10.3.6)]
\lean{Networks.associatedDigraph}
\label{Networks.associatedDigraph}
The \textbf{associated digraph} \texttt{D(G)}: each edge becomes two opposite arcs
(cited as ex 10.3.6).

\emph{The associated digraph \texttt{D(G)}
of a graph \texttt{G} is the digraph obtained when each edge \texttt{e} of \texttt{G} is
replaced by two
oppositely oriented arcs with the same ends as \texttt{e}.}
\end{definition}

\begin{definition}[The vertex-splitting digraph $D'$ of Theorem 11.6, on $V \oplus V$: \texttt{inl v = v'}, \texttt{inr v =\dots}]
\lean{Networks.splitDigraph}
\label{Networks.splitDigraph}
\leanok
The \textbf{vertex-splitting digraph} \texttt{D'} of Theorem 11.6, on \texttt{V $\oplus$
V}: \texttt{inl v = v'}, \texttt{inr v = v''},
with internal arcs \texttt{v' $\to$ v''} and each original arc \texttt{u $\to$ v}
redrawn as \texttt{u'' $\to$ v'}.

\emph{Construct a new
digraph \texttt{D'} from \texttt{D} as follows: (i) split each vertex \texttt{v $\in$ V
\textbackslash{} \{x, y\}} into two
new vertices \texttt{v'} and \texttt{v''}, and join them by an arc \texttt{(v', v'')};
(ii) replace each
arc of \texttt{D} with head \texttt{v $\in$ V \textbackslash{} \{x, y\}} by a new arc
with head \texttt{v'}, and each arc of
\texttt{D} with tail \texttt{v $\in$ V \textbackslash{} \{x, y\}} by a new arc with tail \texttt{v''}.}
\end{definition}

\begin{definition}[IsFeasible]
\lean{Networks.Network.IsFeasible}
\label{Networks.Network.IsFeasible}
\leanok
\uses{Networks.Network, Networks.Network.IsFlow, Networks.Network.fIn, Networks.Network.fOut}
\textbf{Feasibility} for supplies \texttt{$\sigma$} and demands \texttt{dem} (B\&M's
\texttt{$\partial$}).

\emph{Suppose that to each source \texttt{x$_{i}$} of \texttt{N}
is assigned a non-negative integer \texttt{$\sigma$(x$_{i}$)}, called the supply at
\texttt{x$_{i}$}, and to each
sink \texttt{y$_{j}$} of \texttt{N} is assigned a non-negative integer
\texttt{$\partial$(y$_{j}$)}, called the demand at
\texttt{y$_{j}$}.  A flow \texttt{f} in \texttt{N} is said to be feasible if}

f$^{+}$(x$_{i}$) - f$^{-}$(x$_{i}$) $\le$ $\sigma$(x$_{i}$)   for all x$_{i}$ $\in$ X
f$^{-}$(y$_{j}$) - f$^{+}$(y$_{j}$) $\ge$ $\partial$(y$_{j}$) for all y$_{j}$ $\in$ Y

\emph{In other words, a flow \texttt{f} is feasible if the resultant flow out of each
source
\texttt{x$_{i}$} relative to \texttt{f} does not exceed the supply at \texttt{x$_{i}$}, and the resultant flow into
each sink \texttt{y$_{j}$} relative to \texttt{f} is at least as large as the demand at
\texttt{y$_{j}$}.}
\end{definition}

\begin{definition}[$(p,q)$ is realisable by a simple bipartite graph: there is a $(0,1)$-matrix $B$ with row sums $p$ and\dots]
\lean{Networks.RealisableBipartite}
\label{Networks.RealisableBipartite}
\leanok
\texttt{(p,q)} is \textbf{realisable by a simple bipartite graph}: there is a \texttt{(0,1)}-matrix \texttt{B} with row
sums \texttt{p} and column sums \texttt{q}.

\emph{We say that the pair \texttt{(p, q)} is
realisable by a simple bipartite graph if there exists a simple bipartite graph
\texttt{G}
with bipartition \texttt{(\{x$_{1}$, x$_{2}$, \dots{}, x\_m\}, \{y$_{1}$, y$_{2}$,
\dots{}, y\_n\})}, such that
\texttt{d(x$_{i}$) = p$_{i}$} for \texttt{1 $\le$ i $\le$ m} and \texttt{d(y$_{j}$) = q$_{j}$} for \texttt{1 $\le$ j $\le$ n}.}
\end{definition}

\begin{definition}[$G$ is $(m,n)$-orientable: it admits an orientation in which every indegree is $m$ or $n$]
\lean{Networks.IsOrientable}
\label{Networks.IsOrientable}
\leanok
\texttt{G} is \textbf{\texttt{(m,n)}-orientable}: it admits an orientation in which every indegree is \texttt{m} or \texttt{n}.

(B\&M exercise 11.5.5). \emph{An \texttt{(m+n)}-regular graph \texttt{G} is
\texttt{(m,n)}-orientable if it can be oriented so that each indegree is either
\texttt{m} or \texttt{n}.}
\end{definition}

\begin{definition}[$|(S, \bar{S})|$: the number of edges of $G$ crossing the cut $(S, \bar{S})$ (Ex 11.5.5)]
\lean{Networks.edgeBoundaryCard}
\label{Networks.edgeBoundaryCard}
\leanok
\texttt{|[S, $\bar{S}$]|}: the number of edges of \texttt{G} crossing the cut \texttt{(S, $\bar{S}$)} (Ex 11.5.5).

\texttt{[S, $\bar{S}$]} is the set of edges
with one end in \texttt{S} and the other outside.
\end{definition}

\begin{theorem}[Resultant flow out of a set as a sum over singleton vertices]
\lean{Networks.sum_resultant_eq_fOut_sub_fIn}
\label{Networks.sum_resultant_eq_fOut_sub_fIn}
\leanok
\uses{Networks.Network, Networks.Network.fIn, Networks.Network.fOut}
Let $N$ be a network on a finite vertex set $V$, let $f$ assign a nonnegative integer
$f(u,v)$ to each ordered pair of vertices $u,v$, and let $S\subseteq V$. Summing the
resultant flow out of a singleton $\{v\}$ over all vertices $v\in S$ recovers the
resultant flow out of $S$; that is, the following identity holds in $\bbz$:
\[
  \sum_{v\in S}\bigl(f^{+}(\{v\})-f^{-}(\{v\})\bigr)=f^{+}(S)-f^{-}(S).
\]
\end{theorem}

\begin{theorem}[Resultant flow out of the sources equals resultant flow into the sinks]
\lean{Networks.fOut_X_sub_fIn_X_eq_fIn_Y_sub_fOut_Y}
\label{Networks.fOut_X_sub_fIn_X_eq_fIn_Y_sub_fOut_Y}
\leanok
\uses{Networks.Network, Networks.Network.IsFlow, Networks.Network.fIn, Networks.Network.fOut}
Let $N$ be a network with disjoint nonempty source set $X$ and sink set $Y$, and let $f$
be a flow in $N$. For a vertex set $S$, write $f^{+}(S)$ for the total flow along the
arcs directed from $S$ to its complement and $f^{-}(S)$ for the total flow along the
arcs directed from the complement into $S$. Then
\[
f^{+}(X) - f^{-}(X) \;=\; f^{-}(Y) - f^{+}(Y).
\]
\end{theorem}

\begin{theorem}[Value of a flow across an arbitrary cut]
\lean{Networks.val_eq_fOut_sub_fIn}
\label{Networks.val_eq_fOut_sub_fIn}
\leanok
\uses{Networks.Network, Networks.Network.IsCut, Networks.Network.IsFlow, Networks.Network.fIn, Networks.Network.fOut, Networks.Network.val}
Let $N$ be a network with source set $X$, a single distinguished source vertex $x$, and
a single distinguished sink vertex $y$; let $f$ be a flow in $N$; and let $S$ be a set
of vertices forming a cut, so that $x \in S$ and $y \notin S$. Then the value
$\operatorname{val} f = f^{+}(X) - f^{-}(X)$ of the flow is equal to the resultant flow
out of $S$:
\[
\operatorname{val} f = f^{+}(S) - f^{-}(S).
\]
\end{theorem}

\begin{theorem}[Weak duality: flow value is bounded by cut capacity]
\lean{Networks.val_le_capOf}
\label{Networks.val_le_capOf}
\leanok
\uses{Networks.Network, Networks.Network.IsCut, Networks.Network.IsFlow, Networks.Network.capOf, Networks.Network.val}
Let $N$ be a network with disjoint nonempty source set $X$ and sink set $Y$,
distinguished source $x$ and sink $y$, and integer capacity function $c$. Let $f$ be a
flow in $N$, and let $(S,\bar S)$ be a cut, so that $x\in S$ and $y\notin S$. Then the
value of $f$ is at most the capacity of the cut:
\[
\mathrm{val}(f)\;\le\;\operatorname{cap}(S,\bar S)\;=\;\sum_{u\in S}\sum_{v\notin S}
c(u,v).
\]
\end{theorem}

\begin{theorem}[Equality in the flow-value/cut-capacity bound]
\lean{Networks.val_eq_capOf_iff}
\label{Networks.val_eq_capOf_iff}
\leanok
\uses{Networks.Network, Networks.Network.IsCut, Networks.Network.IsFlow, Networks.Network.capOf, Networks.Network.val}
Let $N$ be a network with capacity function $c$, distinguished source $x$ and sink $y$,
let $f$ be a flow in $N$, and let $S$ be a vertex set defining a cut, so that $x \in S$
and $y \notin S$; write $\bar S$ for the complementary vertex set. Then the value of the
flow equals the capacity of the cut, $\mathrm{val}(f) = \mathrm{cap}(S,\bar S)$, if and
only if $f(u,v) = c(u,v)$ for every arc with $u \in S$ and $v \in \bar S$ (each forward
arc across the cut is saturated) and $f(u,v) = 0$ for every arc with $u \in \bar S$ and
$v \in S$ (each backward arc across the cut carries no flow).
\end{theorem}

\begin{theorem}[Equal value and capacity certify a maximum flow and a minimum cut]
\lean{Networks.isMaxFlow_and_isMinCut_of_val_eq_cap}
\label{Networks.isMaxFlow_and_isMinCut_of_val_eq_cap}
\leanok
\uses{Networks.Network, Networks.Network.IsCut, Networks.Network.IsFlow, Networks.Network.IsMaxFlow, Networks.Network.IsMinCut, Networks.Network.capOf, Networks.Network.val}
Let $N$ be a network with single source $x$ and single sink $y$, let $f$ be a flow in
$N$, and let $(S, \bar S)$ be a cut of $N$, so that $x \in S$ and $y \notin S$. If the
value of the flow equals the capacity of the cut, $\mathrm{val}(f) = \mathrm{cap}(S)$,
then $f$ is a maximum flow in $N$ and $(S,\bar S)$ is a minimum cut in $N$.
\end{theorem}

\begin{theorem}[Existence of a maximum flow in every network]
\lean{Networks.exists_maxFlow}
\label{Networks.exists_maxFlow}
\leanok
\uses{Networks.Network, Networks.Network.IsMaxFlow}
Every network $N$ on a finite vertex set admits a maximum flow: there exists a flow $f$
in $N$ such that $\mathrm{val}(f') \le \mathrm{val}(f)$ for every flow $f'$ in $N$.
\end{theorem}

\begin{theorem}[Existence of a minimum cut]
\lean{Networks.exists_minCut}
\label{Networks.exists_minCut}
\leanok
\uses{Networks.Network, Networks.Network.IsMinCut}
Let $N$ be a network with distinguished single source $x$ and single sink $y$. Then $N$
possesses a minimum cut: there exists a set $S$ of vertices with $x \in S$ and $y \notin
S$ whose capacity $\sum_{u \in S}\sum_{v \notin S} c(u,v)$ is least among all sets
separating $x$ from $y$ in this way.
\end{theorem}

\begin{theorem}[Zero maximum flow and minimum cut without a positive-capacity path]
\lean{Networks.maxFlow_and_minCut_zero_of_no_path}
\label{Networks.maxFlow_and_minCut_zero_of_no_path}
\leanok
\uses{Networks.Network, Networks.Network.IsMaxFlow, Networks.Network.IsMinCut, Networks.Network.capOf, Networks.Network.val}
Let $N$ be a network with distinguished single source $x$ and single sink $y$, and
suppose there is no directed path from $x$ to $y$ every arc of which has positive
capacity. Then every maximum flow $f$ in $N$ has value $0$, and every minimum cut $S$ (a
vertex set with $x \in S$ and $y \notin S$) has capacity $0$.
\end{theorem}

\begin{theorem}[Minimum cuts are closed under union and intersection]
\lean{Networks.isMinCut_union_inter}
\label{Networks.isMinCut_union_inter}
\leanok
\uses{Networks.Network, Networks.Network.IsMinCut}
Let $N$ be a network with a single source $x$ and a single sink $y$, and recall that a
cut is determined by a vertex set $S$ with $x \in S$ and $y \notin S$, its capacity
being $\sum_{u \in S}\sum_{v \notin S} c(u,v)$; a cut is a minimum cut if no cut has
strictly smaller capacity. If the cuts determined by vertex sets $S$ and $T$ are both
minimum cuts of $N$, then the cuts determined by $S \cup T$ and by $S \cap T$ are also
minimum cuts of $N$.
\end{theorem}

\begin{theorem}[Max-flow criterion via incrementing paths]
\lean{Networks.maxFlow_iff_no_incrementing_path}
\label{Networks.maxFlow_iff_no_incrementing_path}
\uses{Networks.Network, Networks.Network.IncPath, Networks.Network.IsFlow, Networks.Network.IsMaxFlow}
Let $N$ be a network — a finite digraph carrying a nonnegative integer capacity on each
arc, with disjoint nonempty source and sink sets and a distinguished source $x$ and sink
$y$ — and let $f$ be a flow in $N$. Then $f$ is a maximum flow if and only if there is
no $f$-incrementing path in $N$ ending at the sink $y$.
\end{theorem}

\begin{theorem}[A cut matching the flow when no incrementing path exists]
\lean{Networks.exists_cut_of_no_incrementing_path}
\label{Networks.exists_cut_of_no_incrementing_path}
\leanok
\uses{Networks.Network, Networks.Network.IncPath, Networks.Network.IsCut, Networks.Network.IsFlow, Networks.Network.capOf, Networks.Network.val}
Let $N$ be a network with a single source $x$ and a single sink $y$, and let $f$ be a
flow in $N$. Suppose $N$ contains no $f$-incrementing path from $x$ to $y$ — that is, no
path from the source to the sink each of whose arcs is either a forward arc that is not
yet saturated ($f < c$) or a backward arc carrying positive flow ($f > 0$). Then there
is a cut $(S,\bar S)$, a set of vertices with $x \in S$ and $y \notin S$, whose capacity
equals the value of $f$: \[ \mathrm{val}(f) = \mathrm{cap}(S,\bar S). \]
\end{theorem}

\begin{theorem}[Max-flow min-cut theorem]
\lean{Networks.maxFlow_min_cut}
\label{Networks.maxFlow_min_cut}
\uses{Networks.Network, Networks.Network.IsMaxFlow, Networks.Network.IsMinCut, Networks.Network.capOf, Networks.Network.val}
Let $N$ be a network on a finite vertex set, with capacity function $c$, disjoint
nonempty source and sink sets $X$ and $Y$, and a distinguished source $x$ and sink $y$
(Ford and Fulkerson, 1956). Then $N$ admits a maximum flow $f$ together with a minimum
cut $(S,\bar S)$ — that is, $x \in S$ and $y \notin S$ — whose value and capacity
coincide:
\[
\operatorname{val} f = \operatorname{cap}(S,\bar S).
\]
In particular, the value of a maximum flow in $N$ equals the capacity of a minimum cut.
\end{theorem}

\begin{theorem}[Max-flow value equals arc-disjoint path number in a unit-capacity network]
\lean{Networks.unitCapacity_maxFlow_eq_arcDisjoint}
\label{Networks.unitCapacity_maxFlow_eq_arcDisjoint}
\uses{Networks.Network, Networks.Network.IsMaxFlow, Networks.Network.maxArcDisjointPaths, Networks.Network.val}
Let $N$ be a network with distinguished source $x$ and sink $y$ whose capacity function
takes only the values $0$ and $1$ (that is, every capacity is at most $1$), and let $f$
be a maximum flow in $N$. Then the value of $f$ equals the maximum number of
arc-disjoint directed $(x,y)$-paths in $N$.
\end{theorem}

\begin{theorem}[Min cut equals min arcs destroying all paths in a unit-capacity network]
\lean{Networks.unitCapacity_minCut_eq_minDestroying}
\label{Networks.unitCapacity_minCut_eq_minDestroying}
\uses{Networks.Network, Networks.Network.IsMinCut, Networks.Network.capOf, Networks.Network.minArcsDestroyingPaths}
Let $N$ be a network with capacity function $c$, a single source $x$, and a single sink
$y$, and suppose every arc has capacity at most $1$, that is $c(u,v) \le 1$ for all
vertices $u,v$. For a set $S$ of vertices with $x \in S$ and $y \notin S$, write
$\mathrm{cap}(S) = \sum_{u \in S}\sum_{v \notin S} c(u,v)$ for its cut capacity, and
suppose $S$ is a minimum cut, meaning $\mathrm{cap}(S) \le \mathrm{cap}(S')$ for every
set $S'$ with $x \in S'$ and $y \notin S'$. Then $\mathrm{cap}(S)$ equals the minimum
number of arcs whose deletion destroys every directed $(x,y)$-path in $N$.
\end{theorem}

\begin{theorem}[Menger's theorem: arc version for digraphs]
\lean{Networks.menger_arc_digraph}
\label{Networks.menger_arc_digraph}
\uses{Networks.maxArcDisjointDirectedPaths, Networks.minArcsDestroyingDirectedPaths}
Let $D$ be a finite digraph, and let $x$ and $y$ be two of its vertices. Then the
maximum number of pairwise arc-disjoint directed paths from $x$ to $y$ equals the
minimum number of arcs whose deletion destroys every directed path from $x$ to $y$.
\end{theorem}

\begin{theorem}[Menger's theorem, edge version for graphs]
\lean{Networks.menger_edge_graph}
\label{Networks.menger_edge_graph}
\uses{Networks.maxEdgeDisjointPaths, Networks.minEdgesDestroyingPaths}
Let $G$ be a finite simple graph and let $x$ and $y$ be two of its vertices. Then the
maximum number of pairwise edge-disjoint $(x,y)$-paths in $G$ equals the minimum number
of edges whose deletion destroys every $(x,y)$-path in $G$.
\end{theorem}

\begin{theorem}[Edge connectivity and edge-disjoint paths]
\lean{Networks.k_edge_connected_iff_edge_disjoint_paths}
\label{Networks.k_edge_connected_iff_edge_disjoint_paths}
\uses{Networks.edgeConnectivity, Networks.maxEdgeDisjointPaths}
Let $G$ be a finite simple graph and let $k$ be a natural number. Then the edge
connectivity satisfies $k \le \kappa'(G)$ if and only if, for every pair of distinct
vertices $u$ and $v$, the maximum number of pairwise edge-disjoint paths joining $u$ and
$v$ is at least $k$.
\end{theorem}

\begin{theorem}[Menger's theorem, vertex form for digraphs]
\lean{Networks.menger_vertex_digraph}
\label{Networks.menger_vertex_digraph}
\uses{Networks.maxInternallyDisjointDirectedPaths, Networks.minVerticesDestroyingDirectedPaths}
Let $D$ be a finite digraph and let $x$ and $y$ be two vertices of $D$ with no arc from
$x$ to $y$. Then the maximum number of pairwise internally-disjoint directed paths from
$x$ to $y$ equals the minimum number of vertices (other than $x$ and $y$) whose deletion
destroys every directed path from $x$ to $y$.
\end{theorem}

\begin{theorem}[Menger's theorem, vertex version]
\lean{Networks.menger_vertex_graph}
\label{Networks.menger_vertex_graph}
\uses{Networks.maxInternallyDisjointPaths, Networks.minVerticesDestroyingPaths}
Let $G$ be a finite simple graph, and let $x$ and $y$ be two distinct, nonadjacent
vertices of $G$. Then the maximum number of internally disjoint $(x,y)$-paths in $G$
equals the minimum size of a set of vertices (necessarily distinct from $x$ and $y$)
whose deletion destroys every $(x,y)$-path.
\end{theorem}

\begin{theorem}[Menger's criterion for $k$-connectivity]
\lean{Networks.k_connected_iff_internally_disjoint_paths}
\label{Networks.k_connected_iff_internally_disjoint_paths}
\uses{Networks.maxInternallyDisjointPaths, Networks.vertexConnectivity}
Let $G$ be a simple graph on a finite vertex set $V$, and let $k$ be a natural number
with $|V| \ge k+1$. Then $\kappa(G) \ge k$ if and only if every pair of distinct
vertices $u, v$ of $G$ is joined by at least $k$ internally disjoint $u$–$v$ paths.
\end{theorem}

\begin{theorem}[Dirac's theorem on k vertices lying on a common cycle]
\lean{Networks.dirac_k_vertices_on_cycle}
\label{Networks.dirac_k_vertices_on_cycle}
\uses{Networks.vertexConnectivity}
Let $G$ be a finite simple graph, and let $k \ge 2$ be an integer such that $G$ is
$k$-connected, that is, $k \le \kappa(G)$. Then for any set $T$ of exactly $k$ vertices
of $G$, there is a cycle in $G$ whose vertices include every vertex of $T$.
\end{theorem}

\begin{theorem}[Gale's feasible flow theorem]
\lean{Networks.gale_feasible_flow}
\label{Networks.gale_feasible_flow}
\leanok
\uses{Networks.Network, Networks.Network.IsFeasible, Networks.Network.capOf}
(Gale, 1957.) Let $N$ be a network with source set $X$, sink set $Y$, and capacity
function $c$, and suppose each vertex $x$ is assigned a non-negative integer supply
$\sigma(x)$ and each vertex $y$ a non-negative integer demand $\partial(y)$. Then $N$
admits a feasible flow for the supplies $\sigma$ and demands $\partial$ if and only if,
for every vertex set $S$,
\[
\sum_{v \in Y \cap \bar S} \partial(v) \;-\; \sum_{v \in X \cap \bar S} \sigma(v)
\;\le\; \mathrm{cap}(S, \bar S),
\]
where $\bar S$ denotes the complement of $S$ and $\mathrm{cap}(S,\bar S)$ is the
capacity of the cut from $S$ to $\bar S$.
\end{theorem}

\begin{theorem}[Gale–Ryser condition for bipartite degree sequences]
\lean{Networks.galeRyser_realisable}
\label{Networks.galeRyser_realisable}
\leanok
\uses{Networks.RealisableBipartite}
Let $p=(p_1,\dots,p_m)$ and $q=(q_1,\dots,q_n)$ be sequences of non-negative integers
with equal total, $\sum_{i=1}^{m} p_i = \sum_{j=1}^{n} q_j$, and suppose $q$ is
non-increasing, $q_1 \ge q_2 \ge \cdots \ge q_n$. Then $(p,q)$ is realisable by a simple
bipartite graph if and only if
\[
\sum_{j=1}^{k} q_j \;\le\; \sum_{i=1}^{m} \min(p_i,\, k)
\qquad\text{for every } k = 1, 2, \dots, n.
\]
\end{theorem}

\begin{theorem}[One-row reduction for bipartite degree realisability]
\lean{Networks.realisable_iff_reduced}
\label{Networks.realisable_iff_reduced}
\leanok
\uses{Networks.RealisableBipartite}
Let $p_0, p_1, \ldots, p_m$ and $q_0, q_1, \ldots, q_{n-1}$ be nonnegative integers,
regarded as prescribed row degrees and column degrees. Then the pair $(p, q)$ is
realisable by a simple bipartite graph if and only if the reduced pair $(p', q')$ is
realisable, where $p' = (p_1, \ldots, p_m)$ is obtained by deleting the first row degree
$p_0$, and the reduced column degrees are
\[
q'_j = \max(q_j - 1,\, 0) \quad\text{for every column } j \text{ with } j < p_0, \qquad
q'_j = q_j \quad\text{otherwise;}
\]
that is, one is subtracted (never below zero) from each of the first $p_0$ column
degrees and the remaining column degrees are left unchanged.
\end{theorem}

\begin{theorem}[Cut characterization of $(m,n)$-orientability of regular graphs]
\lean{Networks.isOrientable_iff_exists_partition}
\label{Networks.isOrientable_iff_exists_partition}
\leanok
\uses{Networks.IsOrientable, Networks.edgeBoundaryCard}
Let $G$ be a finite simple graph that is $(m+n)$-regular. Then $G$ is $(m,n)$-orientable
— it admits an orientation in which every vertex has indegree $m$ or $n$ — if and only
if the vertex set can be split into two disjoint sets $V_1$ and $V_2$ whose union is the
whole vertex set, such that for every set $S$ of vertices
\[
\abs{(m-n)\bigl(\abs{V_1\cap S}-\abs{V_2\cap S}\bigr)}\;\le\;\abs{[S,\bar S]},
\]
where $\abs{[S,\bar S]}$ is the number of edges of $G$ with one end in $S$ and the other
outside $S$.
\end{theorem}

\begin{theorem}[Shifting the indegrees of an orientation]
\lean{Networks.isOrientable_pred_succ}
\label{Networks.isOrientable_pred_succ}
\leanok
\uses{Networks.IsOrientable}
Let $G$ be a finite $(m+n)$-regular simple graph, and suppose $n < m$. If $G$ is
$(m,n)$-orientable — that is, its edges admit an orientation in which every vertex has
indegree equal to $m$ or to $n$ — then $G$ is also $(m-1,\,n+1)$-orientable.
\end{theorem}
